\documentclass{article}

\begin{document}

Consider a partially ordered set $P$.

Call a nonempty partially ordered set $D \subseteq P$ {\em countably directed} if every countable subset of $D$ has an upper bound in $D$. 

Call $k \in P$ {\bf $\aleph_1$-compact} if for every countably directed $D \subseteq P$

$$
k\leq\bigvee D\quad\Longrightarrow\quad
k\leq d\mbox{ for some }d\in D .
$$ 


{\bf Lemma.} Consider set $S$ and its powerset $\mathcal P(S)$ partially ordered by $\subseteq$.

Consider a Scott continuous $f:\mathcal P(S)\to\mathcal P(S)$. 

Consider $\mathrm{Fix}(f) = \{x \subseteq S: f(x) = x\}$,
$\mathrm{Fix}(f)$ being partially ordered by $\subseteq$ as well.

Then for every $x\in \mathrm{Fix}(f)$, $x=\bigvee\{k\leq x:k\mbox{ is }\aleph_1\mbox{-compact}\}$.\\

{\bf Remark:} There is {\bf no cardinality restriction on $S$} in this Lemma.\\

{\bf Proof.}

{\bf Fact.} Directed joins in $\mathrm{Fix}(f)$ are unions.

This immediately follows from Scott continuity of $f$:
For every directed $D \subseteq \mathrm{Fix}(f)$, $\bigcup_{d\in D}d = \bigcup_{d\in D} f(d)$, and $f(\bigcup_{d\in D} d) =  \bigcup_{d\in D}f(d)$,
therefore $\bigcup_{d\in D}d$ is a fixed point of $f$, therefore $\bigcup_{d\in D}d \in \mathrm{Fix}(f)$.\\

Now fix $A\in \mathrm{Fix}(f)$ and a token $a\in A \subseteq S$. Scott continuity gives

$$
f(A)=\bigcup_{u\subseteq_{\mathrm{fin}}A}f(u).
$$

Because $A=f(A)$, the formula above means that for each $b\in A$ we can choose a finite witness $u_b\subseteq A$ such that $b\in f(u_b)$.

Starting with $B_0=\{a\}$, define

$$
B_{n+1}=B_n\cup\bigcup_{b\in B_n}u_b,
\qquad
B=\bigcup_{n<\omega}B_n.
$$

Each $B_n$ is finite, so $B$ is countable. ({\bf Note:} countable means either finite or countably infinite.) Every token in $B$ has its chosen witness inside $B$; consequently,

$$
a\in B\subseteq f(B),\qquad B\subseteq A.
$$

We can therefore iterate upwards:

$$
K=\bigcup_{n<\omega}f^n(B).
$$

Scott continuity implies $f(K)=K$. Moreover, $K\subseteq A$, and {\bf $K$ is the least fixed point containing $B$}.

To see that $K$ is $\aleph_1$-compact, suppose $D\subseteq \mathrm{Fix}(f)$ is countably directed and

$$
K\subseteq\bigvee_{\mathrm{Fix}(f)}D=\bigcup D.
$$

The countable set $B$ is contained in the union of countably many members of $D$. Countable directedness supplies one $d\in D$ containing all of $B$. Since $d$ is fixed and  $K$ is the least fixed point containing $B$,

$$
K\subseteq d.
$$

Thus every token of $A$ belongs to an $\aleph_1$-compact fixed point below $A$, proving the Lemma.\\

{\bf Astra gave the following remark:} Notice that {\bf $K$ itself need not be countable}. Its compactness comes from its countable generating set $B$.\\

The reason why $K$ might not be countable is that $f$ is generally allowed to map a finite or countably infinite subset of $S$ to a subset of $S$ which has higher cardinality.


\end{document}